\subsection*{BT990: R3 corrected to the edgewise/fat-tower program}

The R3 continuum origin problem is no longer correctly formulated as a
barycentric-refinement problem.  BT983 showed that the barycentric tower violates
the fatness/shape-regularity hypotheses required by the Cheeger--M\"uller--
Schrader, Dodziuk--Patodi, and FEEC convergence theorems.  The corrected tower is
the edgewise/Freudenthal--Kuhn tower.

The model checks now split the R3 route into two executable halves.  On the
spectral side, BT984 verifies individual Whitney--0/P1 eigenvalue convergence on
the edgewise unit-square tower, and BT985 verifies convergence of the truncated
heat trace
\[
  H_N(t)=\sum_{i\le N} e^{-t\lambda_i}.
\]
At level 6 with 80 eigenvalues, the relative error at \(t=0.05\) is
\(2.1205\times 10^{-3}\).  On the geometric side, BT986 uses projected edgewise
refinement of the octahedral sphere: the minimum angle remains near
\(45^\circ\), the area converges to \(4\pi\), and the Regge scalar integral
\(2\sum_v \delta_v\) agrees with \(8\pi\) to roundoff while local deficits decay.

For the actual curved 4D seeds, BT988 loads the explicit CP\(^2_9\) and
K3\(_{16}\) facets already present in
\texttt{exploration/w33\_explicit\_curved\_4d\_complexes.py}.  Their f-vectors
are
\[
  f(\mathrm{CP}^2_9)=(9,36,84,90,36),\qquad
  f(K3_{16})=(16,120,560,720,288),
\]
with Betti profiles \((1,0,1,0,1)\) and \((1,0,22,0,1)\).  BT989 retires the
barycentric constants \(120/19\) and \(860/19\) for R3 and installs the justified
edgewise top-channel constants:
\[
  \text{top multiplier}=2^4=16,\qquad h_r=2^{-r},\qquad
  \frac{f_4(K3_{16})}{f_4(\mathrm{CP}^2_9)}=rac{288}{36}=8.
\]
The remaining technical task is not to push barycentric refinement harder; it is
to load or derive the local 4-simplex edgewise facet template so the exact lower
incidence and heat-density constants can be recomputed on the CP\(^2_9\)/K3\(_{16}\)
fat tower.
